\section*{v2.4.0 Term-Resolved Channel Specialization Core}
\addcontentsline{toc}{section}{v2.4.0 Term-Resolved Channel Specialization Core}
\begin{quote}
\textbf{Normalization note.} This block is retained as a historical specialization precursor. In the normalized architecture its active content belongs primarily to the projection/translation and field/probe realization layers; it should therefore be read as a later historical expansion, not as a competing core.
\end{quote}

The explicit task of v\docversion\ is to pass from the admissibility language of
v2.3.0 to a first term-resolved structural model. The point is not yet to claim
final physics. Rather, the point is to assign mathematically disciplined roles to
\(T_{\mathrm{tri}},T_{\perp},T_{\mathrm{meas}},T_{\mathrm{res}}\) so that later field equations,
appendix tables, and probe readouts can all be tested against the same channel map.
The guiding principle is that every term must be readable simultaneously in four
ways: as a structural transport rule, as a closure-defect source or sink, as a
probe-interface contribution, and as a future physical interpretation slot.

\subsection*{1. Channel factorization principle}
Let the first-layer master equation be written as
\[
\partial_\tau \Psi
= T_{\mathrm{tri}}(\Psi)
+ T_{\perp}(\Psi)
+ T_{\mathrm{meas}}(\Psi)
+ T_{\mathrm{res}}(\Psi)
- \kappa \nabla\Gamma(\Psi).
\]
In v\docversion\ we impose the factorization rule
\[
T_j = A_j + B_j + C_j,
\qquad j\in\{\mathrm{tri},\perp,\mathrm{meas},\mathrm{res}\},
\]
where:
\begin{itemize}
    \item $A_j$ is the closure-tangent component,
    \item $B_j$ is the explicitly defect-producing normal component,
    \item $C_j$ is a bookkeeping remainder reserved for later model refinement.
\end{itemize}
This factorization is not claimed to be unique. Its function is instead to force
each channel into a common proof grammar: tangent transport, normal defect, and
controlled unresolved remainder.

\paragraph{Definition 1 (channel-adapted splitting).}
Given the projector pair $(P_{\mathrm{tan}},P_{\mathrm{nor}})$ induced by the closure
manifold, define
\[
A_j:=P_{\mathrm{tan}}T_j,
\qquad
B_j:=P_{\mathrm{nor}}T_j,
\qquad
C_j:=0
\]
at the first structural level. More general splittings with nonzero $C_j$ are
allowed when one wants to isolate scale-separated or frame-shifted contributions.
At the present stage, however, the projector-induced splitting is the canonical
choice because it minimizes extra assumptions.

\subsection*{2. Canonical specialization of the four channels}
The structural role of the four channels is now fixed as follows.

\paragraph{Triolic channel.}
The term $T_{\mathrm{tri}}$ is the intrinsic coupling transport generated by the
triolen architecture itself. Its preferred structural property is near tangency:
\[
P_{\mathrm{nor}}T_{\mathrm{tri}} = O(\Gamma).
\]
This expresses that the intrinsic triadic coupling should preserve admissibility to
leading order and create only higher-order closure defects.

\paragraph{Orthogonality channel.}
The term $T_{\perp}$ enforces the standing project constraint that one branch is
perpendicular to the other two. Its preferred structural property is restorative:
\[
\langle \nabla\Gamma, T_{\perp}\rangle \le -\lambda_{\perp}\Gamma
\quad\text{for some }\lambda_{\perp}>0
\]
in the near-closure regime. Thus $T_{\perp}$ is not merely admissible; it is a
coercive correction channel.

\paragraph{Measurement channel.}
The term $T_{\mathrm{meas}}$ encodes the phase-space shift between the measurement
frame and the underlying structural dynamics. It is therefore not expected to be
sign-definite. Instead it is controlled by a relative bound of the form
\[
|\langle \nabla\Gamma, T_{\mathrm{meas}}\rangle|
\le a_{\mathrm{meas}}\|\nabla\Gamma\|^2 + b_{\mathrm{meas}}\Gamma.
\]
This formalizes the standing project rule that measurements do not access the raw
structure directly, but only through a shifted interface.

\paragraph{Residual channel.}
The term $T_{\mathrm{res}}$ is the disciplined storage slot for all contributions not
yet promoted to a stable structural or physical interpretation. In v\docversion\ it is no longer allowed to function as a dumping ground. Instead one requires
\[
T_{\mathrm{res}} = \sum_{m=1}^M T_{\mathrm{res}}^{(m)}
\]
with every summand tagged by its origin class, such as approximation residue,
frame-conversion remainder, truncation error, or phenomenological placeholder.
This turns residuality into an auditable notion.

\subsection*{3. Channel ledger and admissibility matrix}
To make the specialization operational, v2.4.0 introduces the channel
ledger matrix
\[
\Lambda_{\mathrm{ch}}(\Psi)=
\begin{pmatrix}
\ell_{\mathrm{tri}} & \delta_{\mathrm{tri}} & r_{\mathrm{tri}}\\
\ell_{\perp} & \delta_{\perp} & r_{\perp}\\
\ell_{\mathrm{meas}} & \delta_{\mathrm{meas}} & r_{\mathrm{meas}}\\
\ell_{\mathrm{res}} & \delta_{\mathrm{res}} & r_{\mathrm{res}}
\end{pmatrix},
\]
where
\[
\ell_j := \|P_{\mathrm{tan}}T_j\|,
\qquad
\delta_j := \|P_{\mathrm{nor}}T_j\|,
\qquad
r_j := \|C_j\|.
\]
This channel ledger matrix $\Lambda_{\mathrm{ch}}$ is not introduced for numerical decoration. Its function is to give a
single interface where lemma statements, appendix tables, and future simulations
can all report the same structural status.

\paragraph{Proposition 1 (matrix admissibility criterion).}
Suppose there exist constants $\alpha_j,\beta_j\ge0$ such that
\[
\delta_j^2 \le \alpha_j \Gamma + \beta_j \|\nabla\Gamma\|^2
\qquad\text{for each channel }j,
\]
and assume $\sum_j \beta_j < \kappa^2$. Then the channel-specialized master flow is
closure-admissible in a neighbourhood of the closure sector.

\begin{proof}
The normal contributions are exactly the defect-producing pieces entering the
closure balance. By Cauchy--Schwarz,
\[
|\langle \nabla\Gamma, B_j\rangle|
\le \|\nabla\Gamma\|\,\delta_j.
\]
Applying Young's inequality channel by channel and summing yields a bound of the
same type as in the v2.3.0 admissibility criterion, now expressed entirely through
the ledger entries. The compensator term dominates once the cumulative normal
budget stays below the compensator threshold.
\end{proof}

\subsection*{4. Probe-channel specialization}
The same channel grammar must hold for the two-probe sector. Write the probe drift
in the form
\[
\mathcal T_{\mathrm{probe}} =
\mathcal T_{\mathrm{tri}}^{(1)} +
\mathcal T_{\mathrm{tri}}^{(2)} +
\mathcal T_{\perp}^{(Q)} +
\mathcal T_{\mathrm{match}} +
\mathcal T_{\mathrm{res}}^{\mathrm{probe}}.
\]
Here the two intrinsic channels belong to the two probe triolen, the term
$\mathcal T_{\perp}^{(Q)}$ encodes the orthogonal search geometry selected by $Q$,
and $\mathcal T_{\mathrm{match}}$ measures the approach or mismatch between the two
probe traces. The new point of v2.4.0 is that probe matching is promoted to
a channel in its own right rather than being hidden implicitly inside a heuristic
search narrative.

\paragraph{Definition 2 (probe match defect).}
Let
\[
\Delta_{\mathrm{match}}(\Psi_1,\Psi_2)
:= \|\Pi(\Psi_1)-\Pi(\Psi_2)\|^2,
\]
where $\Pi$ is a chosen readout or projection interface. Then the match channel is
called admissible if along the compensated probe flow,
\[
\frac{d}{d\tau}\Delta_{\mathrm{match}}
\le -\eta\,\Delta_{\mathrm{match}} + c\,\Gamma_{\mathrm{probe}}
\]
for some $\eta>0$ and finite $c$.
This means that probe convergence is controlled up to the ambient closure defect.

\paragraph{Corollary 2 (probe convergence under closure control).}
If the probe closure functional remains in a forward-invariant small sublevel tube
and the match channel is admissible, then probe readout mismatch decays up to an
error proportional to the residual closure budget. Hence two probes can be treated
as a mathematically disciplined comparison device rather than only a heuristic
search metaphor.

\subsection*{5. Consequence for appendix tools and future derivations}
The specialization introduced in v2.4.0 changes the status of the appendix
again. From now on, every major table, tool, and shortcut should be readable as a
compressed report of one of the ledger entries or of one of the probe-channel
budgets. This gives a concrete criterion for deciding whether a future tool is
structural or merely expository: a genuine structural tool must reduce, estimate,
or transport one of the channel quantities.

The practical consequence is decisive. Future expansions no longer need to say only
that a section ``supports closure'' or ``helps probe reduction.'' They can specify
which channel is being sharpened, which defect budget is reduced, and which ledger
entry is stabilized. This is exactly the sort of disciplined mathematical
thickening that prepares the transition from structure-first architecture to a more
fully unfolded proof.


